\documentclass{amsart}
\newenvironment{dummy}{}{}

\begin{document}
\begin{dummy}
Without the \texttt{--extra\_mark\_envs} option specified, {\tt corrinline} could not run successfully on this file. 
Though, it shouldn't be hard to just scan for {\tt \string\newenvironment} declarations like is already done for {\tt \string\newtheorem}s.

So, by default, if we have an **ERROUR** in this document, and we mark it in the PDF with an annotation, {\tt corrinline} can't do anything 
because no tex boxes will be created by the marked source since all the text is inside an environment that is not recognized for marking.

Most of the time this isn't an issue because the most frequently used environments which are safe to mark in are already 
tracked by {\tt ALLOWED\_MARK\_ENVIRONMENTS} in {\tt marktex.py}, and frequently-used document-specific ones are typically 
from {\tt \string\newtheorem} declarations, which {\tt corrinline} also finds for the given manuscript source.

Again, it would be very reasonable to track environments from {\tt \string\newenvironment} declarations, and, for good measure, {\tt \string\NewDocumentEnvironment}.

Even more reasonable would be finally just use Sync\TeX{} and not need to worry about the marking at all\dots.
\end{dummy}
\end{document}
